\documentclass{article}
\usepackage{graphicx} % Required for inserting images
\usepackage{amsmath}

\title{\textbf{Hoja de Referencia Prestaciones}}
\author{Juan Diego Bermudez 31.608.369}
\date{May 2026}

\begin{document}

\maketitle

\section{Fórmulas Generales de Prestaciones y Tiempo}
La medida base de las prestaciones es el tiempo de ejecución, y ambas son recíprocas.

\begin{center}
$\text{Prestaciones}_X = \frac{1}{\text{Tiempo de ejecución}_X}$
\end{center}

\begin{center}
$\frac{\text{Prestaciones}_X}{\text{Prestaciones}_Y} = \frac{\text{Tiempo de Ejeción }_X}{\text{Tiempo de Ejecución}_Y}$
\end{center}

\section{Ecuación Clásica de Prestaciones de la CPU}
Esta fórmula desglosa el tiempo de ejecución en tres factores críticos:

Tiempo de ejecución de CPU: \textit{Número de instrucciones × CPI ×Tiempo de ciclo de reloj}

\begin{center}
$\frac{\text{Número de instrucciones×CPI}}{\text{Frecuencia de Reloj}}$
\end{center}

Ciclos de reloj de CPU: Instrucciones del programa×Media de ciclos por instrucción (CPI)

CPI (Ciclos por instrucción):
\begin{center}
CPI = $\frac{\text{Ciclos de reloj de la CPU}}{\text{Número de Instrucciones}}$
\end{center}
IPC (Instrucciones por ciclo): 
\begin{center}
IPC = $\frac{1}{\text{CPI}}$
\end{center}

\section{Tiempo de CPU para un programa}
$$\text{Tiempo de CPU} = \text{Recuento de Instrucciones} \times \text{Tiempo de Ciclo de Reloj}$$

$$\text{Tiempo de CPU} = \frac{\text{Recuento de instrucciones} \times \text{CPI}}{\text{Frecuencia de Reloj}}$$

\section{Rendimiento Relativo}
$$\text{Rendimiento} = \frac{1}{\text{Tiempo de Ejecución}}$$

$$\frac{\text{Rendimiento}_A}{\text{Rendimiento}_B} = \frac{\text{Tiempo de Ejecución}_B}{\text{Tiempo de Ejecución}_A} = n$$

\section{Ciclos de Reloj Totales de CPU}

$$\text{Ciclos del CPU} = \sum_{i=1}^{n} (CPI_i \times \text{Recuento de Instrucciones}_i)$$

\section{CPI Global o Promedio}
$$\text{CPI Global} = \frac{\text{Ciclos de CPU totales}}{\text{Recuento Total de Instrucciones}} = \sum_{i=1}^{n} \left( CPI_i \times \frac{\text{Recuento de Instrucciones}_i}{\text{Recuento Total de Instrucciones}} \right)$$

\section{Ley de Amdahl (Límite del Rendimiento)}
$$\text{Tiempo Nuevo} = \text{Tiempo Viejo} \times \left[ (1 - \text{Fracción Mejorada}) + \frac{\text{Fracción Mejorada}}{\text{Factor de Mejora}} \right]$$


\end{document}
